% Generated draft: requires mathematical and source review.
\documentclass[11pt]{article}
\usepackage[margin=25mm]{geometry}
\usepackage{fontspec}
\setmainfont{Latin Modern Roman}
\usepackage{amsmath,amssymb,mathtools,unicode-math}
\setmathfont{Latin Modern Math}
\usepackage{enumitem,needspace,xurl,hyperref}
\setlength{\parindent}{0pt}
\setlength{\parskip}{5pt}
\setlength{\emergencystretch}{3em}
\newcommand{\N}{\mathbb{N}}
\newcommand{\Z}{\mathbb{Z}}
\newcommand{\Q}{\mathbb{Q}}
\newcommand{\R}{\mathbb{R}}
\newcommand{\C}{\mathbb{C}}
\newcommand{\F}{\mathbb{F}}
\newcommand{\A}{\mathbb{A}}
\newcommand{\PP}{\mathbb{P}}
\newcommand{\E}{\mathbb{E}}
\newcommand{\eps}{\varepsilon}
\newcommand{\dd}{\,\mathrm{d}}
\begin{document}
\section*{UnsolvedMath ID: 7200078}
\textbf{Source problem code:} catalogue entry
\subsection{Definitions and mathematical statement}
Interpret a spherical polyhedron here as a polyhedral region \(P\subset\R^3\) whose boundary is homeomorphic to \(S^2\), not as a polyhedron lying on a round sphere. A translational tiling consists of a locally finite translation set \(\Lambda\subset\R^3\) with
\[
 \bigcup_{\lambda\in\Lambda}(P+\lambda)=\R^3,\qquad
 \operatorname{int}(P+\lambda)\cap\operatorname{int}(P+\mu)=\varnothing\quad(\lambda\ne\mu).
\]
A parallelohedron is a convex polyhedron that tiles Euclidean space by translations.

The source asks whether every nonconvex spherical translational tile \(P\) has its faces groupable into patches with the same combinatorial structure as a parallelohedron. A natural precise reading is whether the facets of \(P\) admit a partition into patches such that merging facets within each patch produces a cell structure on \(\partial P\) whose face incidence poset is isomorphic to that of the boundary of some parallelohedron \(Q\). On that reading the patches represent facets of \(Q\), and the remaining seams and their intersections represent its edges and vertices.

The excerpt does not define patches or what cell structure merging is allowed to produce. In particular, it does not say whether patches must be connected disks, have simple boundary, or consist of coplanar faces. Nor does it explicitly define spherical, state compactness, or require a face-to-face or lattice tiling. Those conventions are therefore unresolved rather than imposed by the displayed tiling definition. The intended claim concerns a combinatorial grouping, not a geometric assertion that \(P\) itself becomes convex or congruent to \(Q\).
\subsection{Short English statement}
Can the faces of every nonconvex translational tile with spherical boundary be merged into patches having the facet, edge, and vertex incidence pattern of a convex parallelohedron? The source leaves the allowed patches and tiling conventions undefined.
\subsection{Sources}
\begin{itemize}
\item Local UnsolvedMath registry, permanent ID 7200078, source code catalogue entry, updated 2026-08-06; no proposer or proposal year supplied.
\item Wikipedia, \emph{List of unsolved problems in mathematics}, Geometry bullet 78, pinned revision 1366636767: \url{https://en.wikipedia.org/w/index.php?title=List_of_unsolved_problems_in_mathematics&oldid=1366636767#Geometry}. Registry access date 2026-07-29; MediaWiki extraction. Wikipedia text is CC BY-SA 4.0; attribution and ShareAlike handling remain required for release. The provenance states unpublished pending primary-source review, despite the true registry published flag.
\item No primary mathematical reference is supplied. The summary explicitly marks status uncertain and reports no verified resolution or citation; its search claims were not independently repeated for this draft.
\end{itemize}
\end{document}
